 \documentclass[12pt,a4paper]{article}
\usepackage[utf8]{inputenc}
\usepackage[greek]{babel} % for greek
%\usepackage[T1]{fontenc}
\usepackage{amsmath}
\usepackage[svgnames]{xcolor}
\usepackage{amsfonts}
\usepackage{amssymb}
\usepackage{graphicx}
\usepackage{hyperref}
\usepackage[left=1.50cm, right=1.50cm, top=3.00cm, bottom=3.00cm]{geometry}
%\author{Stratos Pateloudis}
\title{Ποιοτικά χαρακτηριστικά έξτρα παρθένου ελαιολάδου}
\begin{document}
\maketitle
\section*{Εισαγωγή}
Οι παράγοντες που εππιρεάζουν την ποιότητα του έξτρα παρθένου ελαιολάδου είναι πολλοί.
Γενικώς πέρα απο την οξύτητα σημαντικό ρόλο παίζει η χημική σύσταση του ελαιολάδου όπως η περιεκτικότητά του σε πολυφαινόλες και βιταμίνη Ε, ώστε σε περίπτωση που αυτή είναι υψηλή να χαρακτηριστεί μέχρι και προϊόν υγείας (συμπλήρωμα διατροφής).

Σύμφωνα με τους κανονισμούς της ευρωπαϊκής ένωσης τα επίπεδα αυτά δίνονται στον Πίνακα~\ref{table:health}.
\begin{table}[h!]
	\centering
	\begin{tabular}{|l|c|c|}
		\hline
		\textbf{Παράμετροι} & \textbf{Εξαιρετικό παρθένο} & \textbf{Φαρμακευτικό / Υγείας} \\
		\hline
		Οξύτητα $(FFA)$ & $\leq 0.8\%$ & $< 0.3$–$0.2\%$ \\
		\hline
		Υπεροξείδια & $\leq ~20 ~mEq ~O_2/kg$ & $\leq 5-8 ~mEq~ O_2/kg$ \\
		\hline
		K232 / K270 / $\Delta K$ & $\leq 2.5 / 0.22 / 0.01$ & Ίδια ή μικρότερη \\
		\hline
		Ολικές πολυφαινόλες & $50-1000 ~ mg/kg$ & $\geq ~250 ~mg/kg$ (προτιμότερα $400+ ~mg/kg$) \\
		\hline
		Υδροξυτυροσόλη $\geq$ & – & $5 ~ mg$ ανα $20$ γρ λαδιού \\
		\hline
		Ολεϊκό οξύ & – & $56$–$85\%$ των λιπαρών οξέων \\
		\hline
		Τοκοφερόλες & – & $150$–$250~ mg/kg$ \\
		\hline
		Φυτοστερόλες & – & $100$–$250~ mg$ ανα $100$ γρ \\
		\hline
	\end{tabular}
	\caption{Χαρακτηριστικές τιμές και κριτήρια για εξαιρετικό παρθένο ελαιόλαδο και χαρακτηρισμού ως προϊόντος υγείας.}
	\label{table:health}
\end{table}

\section*{Γενικότητες}
Παραθέτω μερικές γενικότητες για την ποιότητα του ελαιολάδου. Αυτές έχουν να κάνουν με διαφορετικές ποικιλίες, διαφορετικές κλιματολογικές συνθήκες, τρόπος μαζέματος, χρόνος συγκομιδής, αποθήκευση και χρόνος μεταφορά ελαιοκάρπου καθώς και αποθήκευση ελαιολάδου.
\begin{itemize}
	\item Οι διαφορετικές ποικιλίες περιέχουν διαφορετικές περιεκτικότητες σε πολυφαινόλες.   
	\item Οι περιβαλλοντικές συνθήκες εππιρεάζουν την περιεκτικότητα χημικών ενώσεων. Συγκεκριμένα σε θερμές περιοχές (χαμηλό υψόμετρο) τα ελαιόλαδα έχουν μικρότερη περιεκτικότητα ολεϊκού οξέος τα οποία συνδέονται με υψηλότερα επίπεδα παλμικού ή/και λινολεϊκού εξέος\footnote{Κάνοντας μια αναζήτηση, τόσο το ολεϊκό όσο και το λινολεϊκό οξύ, φαίρεται να έχουν ευεργετικές ιδιότητες για την υγεία με το ολεϊκό να είναι πιο ευεργετικό.}. Ωστόσο το τελευταίο μπορεί να είναι σε υψηλή περιεκτικότητα κάτω απο θερμές συνθήκες.   
	\item Η περιεκτικότητα των λιπαρών οξέων της κορονέικης φαίνεται πως δεν αλλάζει. 
	\item Η πολύ βροχή ή/και το πότισμα φαίνεται να μειώνει τις πολυφαινόλες στο ελαιόλαδο, ενώ από την άλλη ο ρόλος των θερμοκρασιών στις πολυφαινόλες φαίνεται να εξαρτάται απο την ποικιλία και δεν είναι απόλυτος.
	\item Φαίνεται ότι πετρώδη χώματα παράγουν ελαιόλαδα με μεγαλυτερη περιεκτικότητα σε πολυφαινόλες καθώς αυτό συνδέεται πιθανόν με την διαθεσιμότητα νερού σε αυτές τις περιοχές. 
	\item Η ωρίμανση των καρπών για συγκομιδή εξαρτάται από την πυκνότητα, δηλαδή από το ΄πόσο φορτωμένο είναι ένα δέντρο. Όσες λιγότερες ελιές έχει ένα δέντρο τόσο πιο γρήγορα ωριμάζουν αυτές, ενώ όσο πιο φορτωμένο είναι τόσο πιο αργά.  Επίσης το χρονικό περιθώριο της <<καλύτερης περιόδου συγκομιδής>> για λίγο φορτωμένα δέντρα είναι μικρότερη από το αντίστοιχο των πολύ φορτωμένων δέντρων. Μια τεχνική που χρησιμοποιείται για την αύξηση των πολυφαινόλων είναι η αποκοπή ελιών, δηλαδή το ανακούφισμα των φορτωμένων δέντρων με μείωση των καρπών κατα 50\% (!) κάτι που συνήθως δεν συνίσταται. Εν τέλει ο φόρτος των δέντρων φαίνεται να μην παίζει σημαντικό ρόλο στην ποιότητα αλλά παίζει σημαντικό ρόλο στον χρόνο συγκομιδής.
	\item Ο χρόνος συγκομιδής εξαρτάται από το φόρτωμα των δέντρων αλλά ως ένας γενικός δείκτης, η καλύτερη περίοδος απο άποψη ποιότητας και ποσότητας ελαιολάδου φαίνεται να είναι όταν ο καρπός έχει περισσότερο πράσινο χρώμα από καφέ, δηλαδή όταν αρχίζει να αλλάζει χρώμα.
    \item Ο τύπος της καλλιέργειας, βιολογική, συμβατική, ολοκληρωμένη, δεν φαίνεται να επηρεάζει την ποιότητα του λαδιού, τόσο στις πολυφαινόλες όσο και στην οξύτητα.  
    \item Σε μερικές ποικιλίες, η φωτοδιαθεσιμότητα φαίνεται πως είναι καταλυτικός παράγοντας για την αύξηση των πολυφαινόλων και μείωση της οξύτητας, ενώ σε άλλες ποικιλίες δεν υπήρξε στατιστική διαφορά. Άρα το κλάδεμα θα πρέπει να γίνεται με τέτοιο τρόπο που να αφήνει το φως να εισέρχεται στους καρπούς, προστατεύοντας ωστόσο τον κορμό από εγκαύματα σε θερμές περιοχές.  
    \item Σε παλιότερες έρευνες εβρέθη ότι το άζωτο συνείσφερε θετικά στην ποιότητα του λαδιού. Σε νεότερες έρευνες ωστόσο, διαπιστώθηκε ότι πολύ υψηλά επίπεδα αζώτου μείωσαν την ποιότητα του λαδιού (περιεκτικότητα σε πολυφαινόλες, οξύτητα, κλπ). Γενικώς η υπεραζώτωση ίσως βοηθάει στην υψηλή παραγωγή καρπού ή/και βλάστησης, όχι όμως τόσο στην ποιότητα λαδιού.
    \item  Η αποθήκευση του καρπού, ο χρόνος αποθήκευσης και η θερμοκρασία παίζουν σημαντικό ρόλο στην ποιότητα του ελαιολάδου. Γενικώς θέλουμε χαμηλές θερμοκρασίες αποθήκευσης και στο ελαιοτριβείο όσο γίνεται πιο γρήγορα με μεταφορά σε αεριζόμενες κάσες και οχι σε τσουβάλια.
    \item Η αποθήκευση του ελαιολάδου είναι επίσης σημαντική και σημαντικό είναι το φιλτράρισμα, καθώς αφαιρεί ένζυμα τα οποία είναι υπεύθυνα για την αλλαγή της γεύσης του καθώς επίσης και για την μείωση των πολυφαινόλων.

\end{itemize}
\section*{Για την ποικιλία χαλκιδικής}
Το ελαιόλαδο από την ποικιλία Χαλκιδικής φαίνεται πως είναι ικανό να πληρεί τα κριτήρια για τον χαρακτηρισμό του ως προϊόν υγείας σύμφωνα με έρευνα του Αριστοτελείου χωρίς κάποια ιδιαίτερη φροντίδα και περιποίηση. 

Σύμφωνα με την έρευνα και τα προηγούμενα στοιχεία, ο χρόνος συγκομιδής, καθώς επίσης και η συλλογή, αποθήκευση και μεταφορά του καρπού παίζουν θεμελιώδη ρόλο στην παραγωγή ενός ποιοτικού, χαμηλής οξύτητας και υψηλής περιεκτικότητας σε πολυφαινόλες, ελαιολάδου. Επίσης η διατήρηση του ελαιολάδου σε σκιερό και ψυχρό χώρο μπορεί να διατηρήσει σχεδόν σταθερά τα παραπάνω επίπεδα (οξύτητα, πολυφαινόλες, κλπ) για τουλάχιστον 18 μήνες. Ένας γενικός κανόνας είναι ότι για την διατήρηση του ελαιολάδου θέλουμε αρκετά χαμηλές θερμοκρασίες (ιδανικά π.χ. $4^{ο}C$) το οποίο σε επαγγελματικές δραστηριότητες επιτυγχάνεται με δεξαμενές ψύξης.

Στην ίδια έρευνα διαπιστώθηκε ότι ο καλύτερος χρόνος συγκομιδής της συγκεκριμένης ποικιλίας είναι όταν βρίσκεται σε επίπεδο ωρίμανσης $\Omega=2$  (δες σχήμα \ref{fig:olive_MI}). Το τελευταίο επίπεδο είναι όταν το πράσινο χρώμα κυριαρχεί σε ποσοστό $\geq 50\%$ στον ελαιόκαρπο όμως μια καλύτερη μέθοδος προσέγγισης είναι ως εξής: γενικά για τις ελιές, η κατηγοριοποίηση της ωρίμανσης του καρπού δίνεται από ένα δείγμα ας πούμε 100 καρπών απο το σύνολο του χωραφιού. Από αυτό το δείγμα οι ελιές χωρίζονται σε 8 κατηγορίες απο το 0 μέχρι το 7 αναλόγως το χρώμα του καρπού. Τότε σε 100 καρπούς ο δείκτης ωρίμανσης $\Omega$ δίνεται ως
\begin{equation}
\Omega =\frac{n_0\cdot 0+n_1\cdot 1+n_2\cdot 2+n_3\cdot 3 +n_4\cdot 4+n_5\cdot 5+n_6\cdot 6+n_7\cdot 7}{100},
\end{equation} 
όπου $n_0,\cdots, n_7$ είναι ο αριθμός των καρπών που βρίσκονται σε αυτή την κατηγορία. Το 0 αφορά τον καρπό που είναι τελείως πράσινος (άγουρος) και το 7 όταν είναι τελείως μαύρος (ζαρωμένος).  Χονδρικά μπορούμε να πούμε ότι το 2 είναι όταν το πράσινο χρώμα κυριαρχεί στον καρπό (δες σχήμα \ref{fig:olive_MI}).

\begin{figure}[ht!]
	\centering 
	\includegraphics[scale=0.5]{olive_MI}
	\caption{Δείκτης ωρίμανσης καρπού}
	\label{fig:olive_MI}
\end{figure}

Εαν δηλαδή απο το χωράφι παρατηρήσουμε 100 τυχαίες ελιές απο τις οποίες 20 είναι στην κατηγορία 1, 30 στην κατηγορία 2 και  50 στην κατηγορία 3 ενώ δεν έχει άλλου χρώματος ελιές τότε ο δείκτης ωρίμανσης είναι 
\begin{equation}
\Omega=\frac{20\cdot 1+30\cdot2+50\cdot3}{100}=\frac{20+60+150}{100}=\frac{230}{100}=2.3,
\end{equation}
που είναι καλή ένδειξη για συγκομιδή. 

Για την μέγιστη περιεκτικότητα ελαιολάδου στον καρπό, πολυφαινόλες και ελάχιστη οξύτητα διαπιστώθηκε ότι ο καλύτερος χρόνος συγκομιδής για την ελιά Χαλκιδικής είναι περίπου όταν $\Omega=2$. Μετά απο αυτό το στάδιο οι πολυφαινόλες αρχίζουν και μειώνονται, η οξύτητα αυξάνει και η ελαιοπεριεκτικότητα μειώνεται. 


Ένας ακόμη σημαντικός παράγοντας ποιότητας ελαιολάδου είναι οι τιμές Κ, οι οποίες εν γένει πρέπει να είναι όσο μικρότερες γίνεται. Χωρίζονται σε δύο κατηγορίες: Κ236 και Κ268 (γνωστές ως 270). 

Η πρώτη κατηγορία (Κ236) ανιχνεύει στερεά σώματα/υπολλείματα στο μήκος κύματος 236$nm$ και η ελαχιστοποίηση αυτού του δείκτη επιτυχγάνεται με την ταχεία μεταφορά των ελαιοκάρπων στον μαλακτήρα. Ο δείκτης αυτός αυξάνεται, εαν οι καρποί αποθηκεύονται σε κακές συνθήκες ή η μεταφορά τους στο ελαιοτριβείο αργήσει.  

Η δεύτερη κατηγορία ανιχνεύει στοιχεία στο μήκος κύματος 268$nm$ και ο δείκτης αυτός αυξάνει όσο το ελαιόλαδο παλαιώνει. Η παλαίωση δεν είναι απαραίτητα μόνο χρονική, καθώς η έκθεση σε ηλιακή ακτινοβολία ή/και υψηλές θερμοκρασίες αυξάνει τον δείκτη Κ270. 

Ενδεικτικά παραθέτουμε έναν πίνακα και ένα γράφημα απο έρευνα που έγινε σε ισπανικά ελαιόλαδα, και τους δείκτες τους. Τα ελαιόλαδα χωρίζονται σε τρεις κατηγορίες: έξτρα παρθένα (extra virgin), παρθένα (virgin), και εξευγενισμένα (Lampante) καθώς οι ενδιαφέρον δείκτες είναι η οξύτητα, ο Κ236 και ο Κ270.    

\begin{table}[ht!]
	\centering
	\begin{tabular}{|c|c|c|c|} 
		\hline
		Δείκτες & Έξτρα παρθένα & Παρθένα & Εξευγενισμένα \\
		\hline 
		\hline
		Οξύτητα	& $0.19$ & $0.34$ & $1.39$ \\ 
		\hline 
		Αιθυλεστέρες λιπαρών οξέων	& $10.7$ & $28.62$ & $ΝΑ$ \\
		\hline
		Κ236	&  $1.64$ & $1.52$ & $ΝΑ$ \\ 
		\hline
		Κ270 &  $0.132$ & $0.133$ & $ΝΑ$\\
		Υπεροξίδεια & $5.62$ & $8.93$ & $13.5$\\
		\hline
	\end{tabular}  
	\caption{Μέσες τιμές δεικτών από 24 λάδια. }
	\label{table:indices}
\end{table}


\begin{figure}[ht!]
	\centering 
	\includegraphics[scale=0.65]{Olive_oil_indices}
	\caption{Δείκτες ελαιολάδων απο τον πίνακα \ref{table:indices}}
	\label{fig:olive_oil_indices}
\end{figure}



\subsection*{Στρατηγική}
Οπότε για να επιτύχουμε εμείς μέγιστες αποδόσεις χωρίς εξωγενείς και μεγάλες ή δύσκολες παρεμβάσεις μπορούμε να κάνουμε τα εξής:
\begin{enumerate}
	\item Συγκομιδή όταν ο δείκτης ωρίμανσης είναι περίπου 2, δηλαδή όταν η ελιά αρχίζει να αλλάζει χρώμα. Αυτό πρέπει να γίνει με απευθείας μέτρηση στο χωράφι με τον παραπάνω τρόπο καθώς όπως αναφέραμε προηγουμένως τα φορτωμένα δέντρα ωριμάζουν πιο αργά τον καρπό απο τα λιγότερο φορτωμένα οπότε η ημερομηνία αλλάζει πάντα. \
	\item Συγκομιδή όσο γίνεται πιο γρήγορα με χρήση ραβδιστικών και γεννήτριας \
	\item Αποθήκευση σε τελάρα που αερίζονται (π.χ. όπως αυτά για ντομάτες και πιστοποιημένα για τρόφιμα) και αποθήκευση σε όσο το δυνατόν χαμηλές θερμοκρασίες\
	\item Ελαχιστοποίηση του χρόνου μεταφοράς μετά την συγκομιδή στο ελαιοτριβείο\
	\item Ίσως και καθημερινή μεταφορά στο ελαιοτριβείο και μετατροπή σε ελαιόλαδο αυθημερόν\ 
	\item  Αποθήκευση ελαιολάδου σε χαμηλές θερμοκρασίες και σκιερό μέρος.
\end{enumerate}
  
\end{document}